\section{The two endpoints}\label{sec:endpoints}

The scaling program is finite by Proposition~\ref{prop:scaling-finite}.
We now identify its endpoint from the pseudo-effectivity of the original
adjoint.  Only ordinary birational comparisons are needed for this last step.

\begin{lemma}\label{lem:pe-step}
Let $(T,B)\dashrightarrow(T',B')$ be an ordinary negative divisorial
contraction or flip between normal compact Kähler klt pairs, with rational
boundaries and rational line-bundle adjoints.  Then $K_T+B$ is pseudo-effective if
and only if $K_{T'}+B'$ is pseudo-effective.
\end{lemma}

\begin{proof}
Set $D=K_T+B$ and $D'=K_{T'}+B'$.  The ordinary comparison in
Proposition~\ref{prop:ordinary-step} and Lemma~\ref{lem:negativity} gives,
on a common smooth compact Kähler resolution $p:V\to T$, $q:V\to T'$,
\[
 p^*D=q^*D'+F,\qquad F\geq0,
\]
where $F$ is an actual divisor exceptional over $T'$.
If $D'$ is pseudo-effective, its pullback plus $[F]$ is pseudo-effective;
descent along $p$ shows that $D$ is pseudo-effective.
Conversely, if $D$ is pseudo-effective, exceptional descent along $q$
applied to $q^*D'+[F]$ proves that $D'$ is pseudo-effective.
Both uses of pullback and descent are justified by
Lemma~\ref{lem:positive-descent}.
\end{proof}

\begin{lemma}\label{lem:mori-not-pe}
Let $g:Y\to S$ be a projective surjective morphism of normal compact Kähler
varieties with $\dim S<\dim Y$.  If a rational line bundle $D$ satisfies
$-D$ $g$-ample, then $D$ is not pseudo-effective.
\end{lemma}

\begin{proof}
Suppose that $c_1(D)$ contains a closed positive current $T$ with local
plurisubharmonic potentials.  Choose a point $y$ where such a potential is
finite and which lies on a positive-dimensional component of a fibre of
$g$.  This is possible on the dense open locus of fibres of the expected
dimension: the pole set of a plurisubharmonic potential contains no open
subset.  That fibre component is projective, so hyperplane sections through
$y$ give an integral curve $C$ through $y$ contained in the fibre.
Let $\nu:C^\nu\to Y$ denote its normalization followed by inclusion.

Local potentials pull back under $\nu$ to subharmonic functions or to the
constant $-\infty$.  Finiteness at a preimage of $y$ excludes the latter
alternative on the connected curve.  Thus they define a positive current
$\nu^*T$ in the class $\nu^*c_1(D)$, and
\[
 D\cdot C=\int_{C^\nu}\nu^*T\geq0.
\]
This contradicts $D\cdot C<0$, which follows from relative ampleness.
The choice of $y$ is essential: no restriction of $T$ to a curve contained
in its pole set has been used.
\end{proof}

\begin{proof}[Proof of Theorem~\ref{thm:finite}]
Run the scaling construction on the given pair $(X,\Delta)$.
Proposition~\ref{prop:ordinary-step} supplies each ordinary negative step,
and Proposition~\ref{prop:scaling-finite} proves that the construction
stops after finitely many such steps.  Write its endpoint as $(Y,\Gamma)$.
All intermediate varieties are normal compact Kähler fourfolds, all
boundaries are the rational transforms of $\Delta$, and all pairs remain
ordinary klt and globally Weil-divisor $\Q$-factorial with the stipulated
canonical datum.  The optional strong global property is preserved by the
same Proposition.  The auxiliary extractions used to prove finiteness do
not change this program, which starts on $X$ itself.

The stopping rule gives either a nef adjoint $K_Y+\Gamma$ or a Mori
contraction $g:Y\to S$.  In the latter case
Proposition~\ref{prop:ordinary-step} supplies a normal compact Kähler base,
a projective surjective morphism with connected fibres, $\dim S<4$,
relative numerical rank $\rho(Y/S)=1$, and relative ampleness of
$-(K_Y+\Gamma)$.  These are precisely the required Mori fibre space
conditions, with relative rank computed from global Cartier classes.
The same Proposition gives relative Bott--Chern dimension one for every
negative contraction, including $g$.

By Lemma~\ref{lem:pe-step}, the adjoint at every stage has the same
pseudo-effectivity status as $K_X+\Delta$.  If the latter is
pseudo-effective, Lemma~\ref{lem:mori-not-pe} excludes the Mori alternative,
so the endpoint is nef.  If it is not pseudo-effective, the endpoint
cannot be nef, since nef classes are pseudo-effective by
Lemma~\ref{lem:positive-descent}; hence the endpoint has the asserted Mori
fibre structure.

For completeness, the nef endpoint has the discrepancy properties of a
log minimal model.  No step extracts a divisor.  On a common resolution of
the finite sequence, with maps $p:V\to X$ and $q:V\to Y$, the ordinary
effective comparison divisors add to
\[
 p^*(K_X+\Delta)=q^*(K_Y+\Gamma)+F,
 \qquad F\geq0,
\]
with $F$ exceptional over $Y$.  Consequently, for every prime divisor $E$
over these models,
\[
 a(E,Y,\Gamma)\geq a(E,X,\Delta).
\]
If a prime divisor on $X$ is contracted, its discrepancy increases strictly
at the divisorial step that contracts its transform and never decreases
thereafter.  Thus the displayed inequality is strict for every such
prime.  Here $a$ denotes the log discrepancy, as throughout the paper.

If the initial adjoint is nef, the construction takes no steps; conversely,
a zero-step nef outcome requires the initial adjoint to be nef.  In the
non-pseudo-effective case a zero-step program is also allowed when the
first selected contraction already gives a Mori fibre structure on $X$.
This completes both endpoint assertions.
\end{proof}
